\section{Budget-normalized causal risk transfer}\label{sec:sharp-transfer}
A stopped task has two different sensitivities. Terminal risk is measured on a probability law, whereas feasibility is measured on the entire live occupation. The preceding resolvent argument respects this distinction. We now use the full mean constraint, rather than its first linear consequence, to obtain a bound with the correct dependence on the acquisition budget. The construction still changes only the routing to an existing terminal label.

\subsection{The exact survival envelope}
For $N\ge1$ put $n=\lfloor N\rfloor$ and $b=N-n$. For $0<\rho\le1$ and $q=1-\rho$, define
\begin{align}
 H_N(\rho)&=\sum_{t=0}^{n-1}q^t+bq^n,
 &\Delta_N(\rho)&=N-H_N(\rho),\label{eq:survival-envelope}\\
 J_N(\rho)&=1-q^{n-1}(1-b\rho).
 \label{eq:replacement-envelope}
\end{align}
The convention $0^0=1$ is used. The temporary symbols $n,b,q$ in this subsection do not denote a deterministic acquisition horizon, a stratum mass, or a likelihood ratio.

\begin{lemma}[Sharp envelopes under a mean constraint]\label{lem:sharp-survival}
Let $\tau$ take values in $\{1,2,\ldots\}$ with $\E\tau\le N$, and let $G$ be independent geometric with success probability $\rho$ on the same set. Then
\begin{equation}\label{eq:sharp-survival-bounds}
 \E\min(\tau,G)\le H_N(\rho),\qquad
 \Pp(G<\tau)\le J_N(\rho)\le\min\{1,\rho(N-1)\}.
\end{equation}
Both first bounds are attained, among all such laws of $\tau$, by putting probabilities $1-b,b$ on $n,n+1$. In addition,
\begin{equation}\label{eq:quadratic-margin}
 \Delta_N(\rho)\ge
 \frac{\rho N(N-1)}{2(1+\rho N)}.
\end{equation}
This last lower bound is not asserted exact.
\end{lemma}
\begin{proof}
For an integer $k\ge1$ the conditional mean of $\min(k,G)$ is
$f(k)=\sum_{t=0}^{k-1}q^t$, and the replacement probability is
$g(k)=1-q^{k-1}$. The successive increments of both functions are nonnegative and nonincreasing. Their piecewise linear interpolants on $[1,\infty)$ are therefore increasing and concave. Jensen's inequality, followed by $\E\tau\le N$, proves the first two assertions, including their equality case. The union bound $1-q^{k-1}\le\rho(k-1)$ proves the remaining bound in \eqref{eq:sharp-survival-bounds}.

For every integer $t\ge0$,
$q^t\le(1+t\rho)^{-1}$: for $0<\rho<1$ this follows from
$(1-\rho)^{-t}\ge1+t\rho/(1-\rho)\ge1+t\rho$, and the endpoints follow directly. Consequently
\[
 \Delta_N(\rho)=\sum_{t=0}^{n-1}(1-q^t)+b(1-q^n)
 \ge \frac{\rho}{1+\rho N}
       \left[\frac{n(n-1)}2+bn\right].
\]
The bracket equals $N(N-1)/2+b(1-b)/2$, proving
\eqref{eq:quadratic-margin}. The equality law in the first assertion is a probability witness, not an uncharged clock in the observer.
\end{proof}

\subsection{A transfer at unchanged resources}
Write $\calR_i^\Theta(N,M)$ for the autonomous value in experiment family $i=1,2$. The two families have the same finite ordered physical space, actions, report reference spaces, legal decoder body and bounded terminal task interface. They may have different preparations, instruments and terminal cost functionals. For each common routing program suppose, uniformly in $\theta$, that the preparation error in base norm is at most $\delta_0$, the induced base-norm error of the joint one-call live/terminal map is at most $\eta$, and the terminal cost functional error per unit positive mass is at most $\alpha$. Integral raw-instrument bounds imply the one-call condition by positivity; no normalized posterior distance is assumed. Every loss belongs to $[0,L]$ in both families.

\begin{theorem}[Sharp-order same-budget transfer]\label{thm:normalized-transfer}
For every real $N\ge1$ and integer $M\ge2$, put
$\zeta=\delta_0+N\eta$ and $c_N=1-N^{-1}$. Then
\begin{equation}\label{eq:normalized-transfer}
 \left|\calR_1^\Theta(N,M)-\calR_2^\Theta(N,M)\right|
 \le \min\left\{L,\ L\left[2\zeta+
          \sqrt{\zeta^2+2c_N\zeta}\right]+\alpha\right\}.
\end{equation}
At $N=1$ the stronger bound $\min\{L,L(\delta_0+\eta)+\alpha\}$ holds. The transfer is effected in each direction by a parameter-independent program with the same state set and the same mean budget. No uniform absorption of the original competitors is required.

The square-root order in $\delta_0$, and in $(N-1)\eta$ for $N\ge2$, is attained up to numerical constants within this experiment class. The witnesses in Theorems~\ref{thm:rare-preparation} and~\ref{thm:rare-instrument} minimize over all two-state programs, not just a chosen regularization family. The constants in \eqref{eq:normalized-transfer} are not claimed best possible.
\end{theorem}
\begin{proof}
Fix any admissible source program. Apply the existing-label regularization of \eqref{eq:killed-routing}. The coupling in Lemma~\ref{lem:margin} and Lemma~\ref{lem:sharp-survival} give source mean at most $H_N(\rho)$ and risk increase at most $LJ_N(\rho)$. The identical regularized program in the other family has stopping floor $\rho$. By Lemma~\ref{lem:resolvent-error}, its mean is at most
\begin{equation}\label{eq:normalized-feasibility}
 H_N(\rho)+\frac{\delta_0+\eta H_N(\rho)}{\rho}
 \le H_N(\rho)+\frac{\zeta}{\rho},
\end{equation}
and its additional risk error is at most $L\zeta+\alpha$. Thus
$\rho\Delta_N(\rho)\ge\zeta$ is a sufficient feasibility certificate.

Assume $N>1$ and $\zeta>0$. Put
\[
 s=\frac{N\zeta}{N-1},\qquad t=s+\sqrt{s^2+2s}.
\]
If $t\le N$, choose $\rho=t/N$. Since $t^2=2s(1+t)$,
\eqref{eq:quadratic-margin} gives $\rho\Delta_N(\rho)\ge\zeta$.
Furthermore
\[
 \rho(N-1)=\zeta+\sqrt{\zeta^2+2c_N\zeta}.
\]
Equations~\eqref{eq:sharp-survival-bounds} and
\eqref{eq:normalized-feasibility} prove the stated one-sided transfer.

If $t>N$, choose $\rho=1$. Both regularized programs stop after exactly one call, so feasibility is exact, independently of any norm estimate. The risk increase is at most $L(N-1)+L(\delta_0+\eta)+\alpha$. The condition $t>N$ is equivalent to
\[
 \zeta>\frac{N(N-1)}{2(N+1)}.
\]
At equality one has $\zeta+\sqrt{\zeta^2+2c_N\zeta}=N-1$; this expression is increasing in $\zeta$. Since $\delta_0+\eta\le\zeta$, the same displayed bound follows.

For $N=1$, every admissible source run already halts after one call. Reroute all formally live outcomes to the fixed terminal fallback. This changes no source execution law, and makes the target mean exactly one; use the one-call norm bound. For $\zeta=0$, the only remaining error is the stated audit error. Take infima over source programs and then exchange the families. Both values lie in $[0,L]$, which justifies the truncation.
\end{proof}

The sharper envelope itself is also an executable certificate. For any $\rho$ satisfying
\begin{equation}\label{eq:exact-envelope-certificate}
 \rho\Delta_N(\rho)\ge\delta_0+\eta H_N(\rho),
\end{equation}
the proof gives error at most
$L[J_N(\rho)+\delta_0+\eta H_N(\rho)]+\alpha$.
A one-call fallback is available even if this sufficient inequality has no solution. In particular, the perturbation enters through the actual paid scale $\delta_0+N\eta$, not through an additional factor $N$ multiplying that scale. The result concerns reoptimized constrained values. It does not say that the unregularized occupation map is continuous, or that one unchanged boundary-optimal program remains feasible after perturbation. This distinction is essential in absorbing control~\cite{PZ,DPR}.

\begin{corollary}[A sufficient finite-precision allocation]\label{cor:normalized-precision}
Suppose $N\ge2$ and $0<a\le1/4$. For any admissible source program choose $\rho=a/N$. If an implementation on the same labels preserves this stopping floor, has preparation and one-call defects satisfying
\begin{equation}\label{eq:precision-allocation}
 \delta_0+N\eta\le a^2/64,
\end{equation}
and terminal cost error at most $\alpha$, it is feasible at the same $N,M$ and has risk at most the source risk plus
$L(a+a^2/64)+\alpha$.
In particular, one-call error of order $a^2/N$ is sufficient, rather than order $a^2/N^2$. This is a sufficient precision region, not an optimal bit-complexity theorem.
\end{corollary}
\begin{proof}
By \eqref{eq:quadratic-margin} and $N-1\ge N/2$,
\[
 \Delta_N(a/N)\ge\frac{a(N-1)}{2(1+a)}\ge\frac{aN}{5}.
\]
The mean perturbation in \eqref{eq:normalized-feasibility} is at most $aN/64$, strictly less than this margin. The risk assertion follows from $J_N(a/N)\le a$ and Lemma~\ref{lem:resolvent-error}. Report averaging and rounding before regularization, as in Lemma~\ref{lem:finite-routing}, preserve the floor. For $1<N<2$ one may instead use \eqref{eq:exact-envelope-certificate}; no uniform small-margin assertion is being omitted at $N=1$.
\end{proof}
